\scanpage{28}

Donc on trouve, \deleted{pour} $S$ fixé, une application
\[
 \mathrm{Pl}_{\mathrm{isot}}(X,S)\longrightarrow
 \begin{gathered}\text{Ensemble des systèmes}\\
 \alpha,\pi,g,\nu\ \text{associés à }X\quad(\text{à expliciter}\\
 \text{combinatoirement})\end{gathered}
\]
et on est donc ramené à étudier, pour $\rho=\alpha,\pi,g,\nu$ fixé l'ens des $\mathrm{Pl}_{\mathrm{isot}}(X,S;\rho)$ qui correspondent à ce $\rho$ (non vide sss $S\simeq S_\rho\overset{\text{déf}}=S_{\alpha,\pi,g,\nu}$). \margin{ce qu'on va supposer dorénavant} Or on a une application \deleted{\unsure{Isotopt|Isot}} ``restriction à $X$
\[
 \boxed{\mathrm{Isot}(S_\rho,S)\xrightarrow{r}\mathrm{Pl}_{\mathrm{isot}}(X,S;\rho)}
\]
\inserted{\unsure{préservant|et précisé} orientation\ldots.} qui, d'après ce qu'on vient de voir, est \uline{surjective}.

\deleted{Il reste à}

\deleted{b) étudier quelles sont les fibres de cette application, \deleted{i.e.}

a) Déterminer l'ens $\mathrm{Isot}(S_\rho,S)$

La question a) est classique et résolue : $\mathrm{Isot}(S_\rho,S)$ est un torseur \inserted{à dr} sous $\mathrm{Autext}(\pi_1(S_\rho))$ \inserted{(ou à g. sous $\mathrm{Autext}(\pi_1(S))$)} [du moins si $S_\rho$ est compacte i.e. $\sum_{\lambda\in\Lambda}\unsure{\nu|\lambda}_\lambda=0$ i.e. les $\nu_\lambda$ nuls \ldots]. \deleted{Donc la question est la suivante}

Soit $f:S_\rho\xrightarrow{\sim} S_\rho$ un homéomorphisme, dont la classe d'homotopie est déterminée par l'\deleted{\unsure{endom}}\inserted{extérieur} induit par $f$ sur $\pi_1(S_\rho)=G_\rho$, \deleted{soit \unsure{$E(f)$}} \inserted{soit $E(f)\in\mathrm{Autext}(G_\rho)$}, considérons $f|X=f_X$. Quelle condition sur $E(f),E(g)$ assurent que $f_X\overset{\mathrm{isot}}{\simeq}g_X$ ?}

\diagram{\begin{tikzcd}[ampersand replacement=\&,row sep=small,column sep=small]
G_\rho\arrow[r,shift left=.55ex,"{\bar f,\bar g}"]\arrow[r,shift right=.55ex]\&G_\rho\\
\pi_1(X)\arrow[u]\arrow[ur,bend right=48,dash]\&
\end{tikzcd}}
\uline{NB} Notons que le premier
\[
 G=\mathrm{Isot}_{\mathrm{or}}(S,S)\qquad[\simeq G_\rho\simeq\mathrm{Isot}_{\mathrm{or}}(S_\rho,S_\rho)]
\]
opère sur les deux mbres de $(*)$, et que $r$ commute à l'opération de ce groupe. D'autre part, \deleted{on sait}\inserted{il est clair} que le premier mb est un torseur sous $G$, donc, choisissant $f:S_\alpha\simeq S$, \unsure{et si|si} $r(f)=f_X=f|X$ est son image, et $H$ son stabilisateur, on a que $r$ s'identifie à l'application $G\longrightarrow G/H$. Donc on
% archive sheet number: 27
